\section{Measuring with a language model}
\label{sec:llm-measurement}

This section exists so the literature specimen has somewhere to land. Every key
below is claimed on \texttt{QBt4-for-literature.md} and resolves in the shared
bank, and the sentences say only what those sources actually say.

Large language models acquired their general few-shot behaviour at scale, which
is what made them usable as measurement instruments without task-specific
training \citep{brown2020language}. Applied to personality specifically, a
language model has been used to augment text before a smaller model performs
the detection, rather than being trusted to produce the label alone
\citep{hu2024llm}. Surveys of medical applications catalogue the same
opportunity alongside the same reservations about reliability and evaluation
\citep{Nazi2024LLMHealthcare}. None of this settles whether an online rating is
a valid signal about a physician in the first place, which is a separate
question with its own evidence: patient ratings have been compared directly
against physician peer review \citep{mcgrath2018validity}.

The point for a page-type specimen is mechanical rather than substantive. Each
of the four keys was CLAIMED on a literature page with a sentence saying what
job the source does, and each is cited here. The bibliography this document
prints was generated from those two facts and from nothing else: no key was
copied into a local file by hand, and a key that no page claims and no section
cites does not appear.
